\documentclass[12pt]{ctexart}
\author{Yan}
\usepackage{amsmath,amssymb}
\usepackage{geometry}
\usepackage{setspace}
\usepackage{multirow}
\usepackage{array}
\usepackage{makecell}
\geometry{a4paper, margin=1in}

\geometry{a4paper,margin=2.5cm}
\setstretch{1.25}

\title{拍拍Clapclap 双人版 规则书}
\date{2026.3.17}

\begin{document}

\maketitle
\tableofcontents
\newpage

%------------------------------------------------
\section{游戏介绍}

拍拍（Clapclap）是一款多人参与的同步出招战斗游戏，所有玩家围成一圈，通过击掌与手势的组合进行回合制对抗。

在游戏中，玩家需要在资源积累、攻击、防御与锦囊反制之间进行权衡，通过对对手行动的判断与博弈逐步取得优势，最终成为唯一存活的玩家。

与传统回合制游戏不同，本游戏的所有行动均在同一时刻完成，这意味着玩家无法依赖先后手优势，
而必须通过对局势的判断与对对手心理的预判来取得胜利。

本规则书将详细介绍双人版游戏的基本规则、回合流程、资源系统、攻击与防御机制以及锦囊反制等核心玩法，
帮助玩家快速上手并享受游戏的乐趣。

%------------------------------------------------
\section{游戏开始以及资源介绍}

游戏开始时，每位玩家都拥有相同的初始状态。

这个游戏总共有六种资源，分别是：生命值、气、盾、火种、电池与镐。

每位玩家初始拥有若干点生命值，一般为1点，其余资源均为0。

\begin{itemize}
    \item 生命值：代表玩家的生存状态，当生命值降至0时，该玩家立即出局。
    \item 气：基础资源之一，可以通过特定手势积累。
    \item 盾：基础资源之一，可以通过特定手势积累。
    \item 火种：由盾系攻击衍生出的中间资源。
    \item 电池：由盾系攻击衍生出的中间资源。
    \item 镐：特殊资源，可以抵挡一点伤害。
\end{itemize}

游戏中除镐外，其余资源均没有上限。

%------------------------------------------------
\section{游戏进行}

所有玩家在游戏中需要保持固定的姿势：右手手掌向下，左手手掌向上，使得每个人可以同时与左右两侧的玩家进行击掌。

游戏按回合进行，每一回合分为三个阶段。

\begin{itemize}
    \item 击掌阶段：所有玩家同时与左右两侧各击掌一次，形成一个闭环的互动。
    \item 出手阶段：所有玩家必须同时做出一个手势。
    \item 结算阶段：根据所有玩家的手势进行统一判定，包括攻击、防御、锦囊效果以及资源变化。
    血量为0及以下的玩家视为死亡，被移除游戏
\end{itemize}

直到最后只剩下一名玩家存活时，游戏结束，该玩家获胜。

出手阶段不能延迟或观察他人后再决定，慢出手的玩家被称作“蛤蟆”，出不合规手势的玩家被称作“蟆蛤”，两者都将直接死亡。

%------------------------------------------------
\section{手势介绍}

游戏中的所有手势可以分为四大类：资源、攻击、防御和锦囊。

每个手势都有其攻击力和防御力。

攻击与防御的核心判定规则是：当一方大于另一方的防御力时，攻击成立并造成伤害；双方攻击力相等时，双方视为对掉。

\subsection{资源手势}

资源手势包括气和盾。

\begin{itemize}
    \item 气：获得1格气，防御力为0
    \item 盾：获得1格盾，防御力为1.5
\end{itemize}

资源手势的攻击力为0，但可以积累资源，为后续的攻击和防御提供基础。

\subsection{攻击手势}

攻击手势分为气系攻击和盾系攻击。所有攻击手势在具备攻击力的同时，其防御力也等于其攻击力。

\begin{itemize}
    \item gi：消耗1格气，攻击力为1
    \item 破：消耗2格气，攻击力为2
    \item 冷锋：消耗3格气，攻击力为3
    \item 如来：消耗5格气，攻击力为4
    \item 黑洞：消耗8格气，攻击力为5
    \item Fire：消耗2格盾，攻击力为1.5，并获得1个火种
    \item 闪电：消耗3格盾，攻击力为2，并获得1个电池
    \item 烈焰：消耗4格盾或2个火种，攻击力为3
    \item Shining：消耗6格盾或2个电池，攻击力为4
\end{itemize}

除去如来和Shining可以造成两点伤害、黑洞可以造成三点伤害以外，其它攻击都只能造成一点伤害。

Shining可以拆分为两个闪电分别结算；
黑洞可以拆分为三个小黑洞，攻击力仍为5，但每段仅造成1点伤害，可以分配给一个或多个目标。

\subsection{防御手势}

防御手势用于抵御敌方攻击，攻击力为0。

\begin{itemize}
    \item 十字：消耗2格气，防御力为3
    \item 八卦：消耗3格气，防御力为4
\end{itemize}

这些防御手势可以有效抵挡中高强度攻击，是对抗气系和盾系爆发的重要手段。

\subsection{锦囊手势}

锦囊手势为游戏提供了关键的反制与博弈机制。

\begin{itemize}
    \item 你吃：消耗1格气，可以针对特定攻击进行克制。
    若目标为闪电，则其攻击失效且无法获得电池，若目标为破，则其使用者会受到1点反噬伤害。
    \item 双吃：消耗2格气，具备你吃的全部能力，还可以以Shining为目标使其攻击失效且无法获得电池。
    同时也可以拆分为两个你吃，分别针对不同目标。
    \item 闪：无消耗，每局最多使用两次，使用后在本回合退出游戏，不参与任何结算。
    \item 镐：消耗2格气，使用后获得一个镐。
\end{itemize}

你吃和双吃未命中目标时，其消耗依然生效，同时防御力视为0。

镐具有独特的风险机制。当一名玩家持有两个或以上的镐时，会立即触发爆镐，直接死亡；
gi可以在对方使用镐的当回合“抢镐”，对方不仅无法获得镐，自己还会多一个镐。
因此镐既是恢复手段，也是潜在危险资源。

%------------------------------------------------
\section{总结}

拍拍（Clapclap）在简单的操作形式之下构建了一个高度紧凑的战斗系统。
玩家既需要通过资源积累建立优势，又必须在关键时刻做出正确判断。
气体系提供了稳定而高效的战斗能力，盾体系则带来多样化的发展路径与资源转化方式，
而锦囊机制进一步强化了针对性与不确定性，使得每一回合都充满变数。
游戏的胜负往往取决于关键时刻的选择，而非单纯的资源数量，这也使得对局始终保持紧张与动态的平衡。

在实际体验中，拍拍（Clapclap）既保留了拍手互动带来的节奏感与参与感，又通过复杂而统一的规则体系赋予其足够的策略深度。
它既可以作为轻量的互动游戏快速展开，也可以在熟练玩家之间演变为高强度的心理博弈。
这种简单形式与深层策略并存的特性，正是其最核心的魅力所在。

%------------------------------------------------
\newgeometry{left=1cm,right=1cm,top=2.5cm,bottom=2.5cm}
\section*{附录：手势一览表}

\renewcommand\arraystretch{1.3}

\centering
\begin{tabular}{|>{\centering\arraybackslash}m{1.2cm}|c|c|c|c|c|}
\hline
\textbf{类型} & \textbf{名称} & \textbf{消耗} & \textbf{攻击力} & \textbf{效果及备注} & \textbf{手势} \\
\hline

\multirow[c]{2}{*}{\centering\makecell{资\\源}} & 气 & & 0 & 获得1格气 & 双手作鼓掌状 \\ \cline{2-6}
 & 盾 & & 1.5 & 获得1格盾 & 大臂紧贴身体，双拳握于胸前\\ \cline{2-5}

\hline

\multirow[c]{8}{*}{\centering\makecell{攻\\击}} & gi & 1格气 & 1 & & \makecell{大臂紧贴身体，小臂与大臂垂\\直，双手握拳向前} \\ \cline{2-6}
 & 破 & 2格气 & 2 & & 双手手掌朝前，向前推出 \\ \cline{2-6}
 & 冷锋 & 3格气 & 3 & & 右手三根手指作手枪状 \\ \cline{2-6}
 & 如来 & 5格气 & 4 & &  \makecell{右手手掌向左，除食指弯曲之\\外其它手指竖直向上}\\ \cline{2-6}
 & 黑洞 & 8格气 & 5 & &  \makecell{左手小臂平行置于右手小臂上\\方，双手指尖置于肘前区形成\\朝前方的矩形}\\ \cline{2-6}
 & Fire & 2格盾 & 1.5 & 获得1个火种 & \makecell{双手交叉于腹前，手掌向上\\手指自然弯曲}\\ \cline{2-6}
 & 闪电 & 3格盾 & 2 & 获得1个电池 & \makecell{双手置于身前，手心向下，伸\\出食指和中指并拢，用右手双\\指第二指关节敲击左手双指第\\二指关节}\\ \cline{2-6}
 & 烈焰 & \makecell{4格盾 或\\2个火种} & 3 & & \makecell{双手交叉于腹前，手掌向下，\\手指自然弯曲} \\ \cline{2-6}
 & Shining & \makecell{6格盾 或\\2个电池} & 4 & & \makecell{双手各三根手指作手枪状置于\\额头前上方区域，手背向外，\\用右手双指第二指关节敲击左\\手双指第二指关节} \\

\hline
\multirow[c]{2}{*}{\centering\makecell{防\\御}} & 十字防 & 2格气 & 3 & & \makecell{双手握拳置于身前，左手小臂\\平行置于右手小臂上方}\\ \cline{2-6}
 & 八卦 & 3格气 & 4 & & \makecell{双手小臂置于身前上下移动，\\手掌面向自己，手指自由弯曲} \\

\hline
\multirow[c]{4}{*}{\centering\makecell{锦\\囊}} & 你吃 & 1格气 & 0 & 吃掉闪电吃死破 & 伸出食指指向前方 \\ \cline{2-6}
 & 双吃 & 2格气 & 0 & \makecell{吃掉闪电吃死破\\吃掉Shining} & 伸出食指和中指并拢指向前方 \\ \cline{2-6}
 & 闪 & 无 & 无 & 每局最多使用两次 & 伸出一根手指在身前晃动 \\ \cline{2-6}
 & 镐 & 2格气 & 无 & 获得1个镐 & 右手握拳锤向心脏位置\\

\hline

\end{tabular}

\restoregeometry

\end{document}